\documentclass[12pt,a4paper]{article}
\usepackage[utf8]{inputenc}
\usepackage[bulgarian]{babel}
\usepackage[T2A]{fontenc}
\usepackage{graphicx}
\usepackage{amsmath}
\usepackage{amsfonts}
\usepackage{amssymb}
\usepackage{listings}
\usepackage{listingsutf8}
\usepackage[dvipsnames,table]{xcolor}
\usepackage{geometry}
\usepackage{hyperref}

% Define custom colors
\definecolor{lightblue}{RGB}{173,216,230}
\definecolor{lightgreen}{RGB}{144,238,144}
\definecolor{lightyellow}{RGB}{255,255,224}
\definecolor{lightcoral}{RGB}{240,128,128}
\usepackage{tikz}
\usepackage{tikz-uml}
\usepackage{float}
\usepackage{array}
\usepackage{booktabs}
\usepackage{fancyhdr}
\usepackage{subcaption}

\geometry{
    left=2.5cm,
    right=2.5cm,
    top=3cm,
    bottom=3cm
}

\pagestyle{fancy}
\fancyhf{}
\rhead{MetroSim - Документация}
\lhead{\leftmark}
\cfoot{\thepage}

% Настройки за кода
\lstset{
    basicstyle=\ttfamily\footnotesize,
    keywordstyle=\color{blue}\bfseries,
    commentstyle=\color{gray},
    stringstyle=\color{red},
    numberstyle=\tiny\color{gray},
    numbers=left,
    frame=single,
    breaklines=true,
    showstringspaces=false,
    language=Java,
    captionpos=b,
    inputencoding=utf8,
    extendedchars=true
}

\hypersetup{
    colorlinks=true,
    linkcolor=blue,
    urlcolor=blue,
    pdftitle={MetroSim - Документация},
    pdfauthor={Александър Цветанов}
}

\title{\Huge\textbf{MetroSim} \\
\Large Симулатор на транспортна мрежа с паралелна обработка \\
\large Документация на проекта}

\author{Александър Цветанов}
\date{\today}

\begin{document}

\maketitle
\newpage

\tableofcontents
\newpage

\section{Въведение}

\subsection{Обща информация за проекта}

\textbf{MetroSim} е десктоп симулатор на транспортна мрежа, който визуализира движението на превозни средства в опростена метро/автобусна система. Проектът е изграден с \textbf{Java 17} и използва \textbf{Swing} за графичния интерфейс.

Основната цел на приложението е да демонстрира как множество \textbf{шаблони за проектиране (Design Patterns)} могат да бъдат комбинирани в \textbf{паралелно, събитийно и разширяемо приложение}. Това го прави отлична примерна имплементация за университетски курсове по:

\begin{itemize}
    \item \textbf{Шаблони за проектиране} – демонстрация на практическото приложение на GoF шаблоните
    \item \textbf{Паралелно програмиране} – безопасна работа с нишки в графични приложения
    \item \textbf{Софтуерна архитектура} – изграждане на модулни и разширяеми системи
\end{itemize}

\subsection{Ключови характеристики}

\begin{itemize}
    \item \textbf{Реално време симулация} на превозни средства и световен часовник
    \item \textbf{Отзивчив Swing интерфейс} – актуализации чрез неизменяеми \texttt{WorldSnapshot} обекти
    \item \textbf{Безопасен паралелизъм} използвайки \texttt{ForkJoinPool} и \texttt{ScheduledExecutorService}
    \item \textbf{Разширяем дизайн} – нови маршрутизатори, градове и метрики чрез плъгини
    \item \textbf{Command pattern каркас} за бъдещи функции за редактиране/undo-redo

\end{itemize}

\subsection{Технически изисквания}

\begin{itemize}
    \item Java Development Kit (JDK) 17 или по-нова версия
    \item Gradle за управление на зависимостите и билд процеса
    \item Операционна система: Windows, macOS, Linux с поддръжка на Java Swing
\end{itemize}

\section{Идея и концепция на приложението}

\subsection{Основна идея}

Основната идея зад MetroSim е да се създаде симулатор, който моделира реалистично поведение на транспортна мрежа. Всяко превозно средство в системата:

\begin{enumerate}
    \item \textbf{Започва} от определен възел (позиция) в мрежата
    \item \textbf{Движи се} стъпка по стъпка към крайната си дестинация
    \item \textbf{Спира} за кратък престой („dwell") при достигане на дестинацията
    \item \textbf{Обръща посоката} и продължава към нова дестинация
    \item \textbf{Работи едновременно} с всички останали превозни средства в системата
\end{enumerate}

\subsection{Симулационен модел}

Симулацията се основава на \textbf{дискретно време} с фиксирани интервали (tick-ове). На всеки tick всички превозни средства актуализират своето състояние паралелно. Това позволява:

\begin{itemize}
    \item \textbf{Детерминистично поведение} – резултатите са предвидими и възпроизводими
    \item \textbf{Ефективна паралелизация} – всяко превозно средство може да се обработва независимо
    \item \textbf{Лесно тестване} – състоянието може да се проверява на всеки времеви интервал
\end{itemize}

\subsection{Визуализация и взаимодействие}

Потребителският интерфейс предоставя:

\begin{itemize}
    \item \textbf{Карта} – визуализация на движещите се превозни средства
    \item \textbf{Контролни бутони} – Start/Pause за управление на симулацията
    \item \textbf{Панел за инспекция} – детайлна информация за всяко превозно средство
    \item \textbf{Времева линия} – индикатор за напредъка на симулацията
\end{itemize}

\section{Архитектура на системата}

\subsection{Слоеста архитектура}

MetroSim използва слоеста архитектура с ясно разделение на отговорностите:

\begin{table}[H]
\centering
\begin{tabular}{|p{3cm}|p{10cm}|}
\hline
\textbf{Слой} & \textbf{Описание} \\
\hline
\texttt{app/} & Главна точка на стартиране и инициализация на компонентите \\
\hline
\texttt{ui/} & Графичен интерфейс (Swing) – прозорци, панели и медиатор \\
\hline
\texttt{domain/} & Основна симулационна логика – свят, превозни средства, състояния \\
\hline
\texttt{infrastructure/} & Конкурентен двигател и фасада за управление на симулацията \\
\hline
\texttt{services/} & Алгоритми за маршрутизиране, зареждане на сценарии и плъгини \\
\hline
\texttt{metrics/} & Система за измервания и статистики \\
\hline
\texttt{themes/} & Теми и цветови пакети чрез плъгини \\
\hline
\texttt{events/} & Събитийна система (Observer) за обновяване на интерфейса \\
\hline
\end{tabular}
\caption{Слоеве в архитектурата на MetroSim}
\end{table}

\subsection{Модел на паралелност}

Системата използва няколко типа нишки за различни задачи:

\begin{table}[H]
\centering
\begin{tabular}{|p{4cm}|p{9cm}|}
\hline
\textbf{Тип нишка} & \textbf{Отговорност} \\
\hline
EDT (Event Dispatch Thread) & Обработка на графиката и взаимодействие с потребителя \\
\hline
Симулационен планировчик & Задейства цикъла на симулацията на определени интервали \\
\hline
Работен пул (ForkJoinPool) & Изпълнява обновленията на превозните средства паралелно \\
\hline
\end{tabular}
\caption{Модел на паралелност в MetroSim}
\end{table}

\section{Детайлен анализ на класовете}

\subsection{Domain слой}

\subsubsection{Класове и интерфейси}

\paragraph{World}
Централният клас, който съхранява състоянието на цялата симулация:

\begin{lstlisting}[caption=Основни методи на класа World]
public final class World {
    private final List<Vehicle> vehicles = new ArrayList<>();
    private final SimClock clock = new SimClock();

    public List<Vehicle> vehicles() { return vehicles; }
    public SimClock clock() { return clock; }
    public void advanceClock() { clock.advance(1); }
    public WorldSnapshot snapshot(long tickId) { ... }
}
\end{lstlisting}

\paragraph{Vehicle}
Представлява превозно средство с вътрешни състояния:

\begin{itemize}
    \item \textbf{Idle} – започално състояние
    \item \textbf{EnRoute} – движение към дестинация
    \item \textbf{Dwell} – престой на спирка
\end{itemize}

\paragraph{VehicleState}
Интерфейс за имплементация на State pattern:

\begin{lstlisting}[caption=Интерфейс VehicleState]
public interface VehicleState {
    void onTick(Vehicle ctx, SimClock clock, CancellationToken ct);
    String name();
}
\end{lstlisting}

\paragraph{WorldSnapshot и VehicleSnapshot}
Неизменяеми обекти (records), които представляват състоянието в даден момент:

\begin{lstlisting}[caption=Snapshot класове]
public record WorldSnapshot(long tickId, List<VehicleSnapshot> vehicles) {
    public static final WorldSnapshot EMPTY = new WorldSnapshot(0, List.of());
}

public record VehicleSnapshot(String label, int nodeIndex, String stateName) {}
\end{lstlisting}

\subsection{Infrastructure слой}

\subsubsection{SimulationEngine}

Ядрото на симулацията, което обработва всички превозни средства паралелно:

\begin{lstlisting}[caption=Ключов метод на SimulationEngine]
public WorldSnapshot step(CancellationToken ct) {
    List<Callable<Void>> tasks = new ArrayList<>();
    for (Vehicle v : world.vehicles())
        tasks.add(() -> {
            v.tick(world.clock(), ct);
            return null;
        });
    invokeAllCancellable(tasks, ct);
    world.advanceClock();
    return world.snapshot(tick.incrementAndGet());
}
\end{lstlisting}

\subsubsection{SimulationFacade}

Фасада, която скрива сложността на управлението на нишки и синхронизация:

\begin{lstlisting}[caption=SimulationFacade управление]
public synchronized void start() {
    if (token.isCancelled()) {
        token = new CancellationToken();
    }
    runner = scheduler.scheduleAtFixedRate(
        () -> {
            if (token.isCancelled()) return;
            WorldSnapshot s = engine.step(token);
            snapshots.publish(s);
        }, 0, 50, TimeUnit.MILLISECONDS);
}
\end{lstlisting}

\subsection{Services слой}

\subsubsection{Routing подсистема}

Системата за маршрутизиране използва Strategy pattern с няколко имплементации:

\begin{itemize}
    \item \textbf{TimeOptimalRouter} – базов алгоритъм за оптимално време
    \item \textbf{ThirdPartyRouterAdapter} – адаптира външни библиотеки
    \item \textbf{LoggingRouter} – декоратор за логване
    \item \textbf{MeteredRouter} – декоратор за метрики
\end{itemize}

\subsection{UI слой}

\subsubsection{MainWindow}

Главният прозорец на приложението, който координира всички панели:

\begin{lstlisting}[caption=Структура на MainWindow]
public MainWindow(Bootstrap.Components c) {
    // Създаване на панели
    MapPanel map = new MapPanel(c.cityPack);
    TimelinePanel timeline = new TimelinePanel();
    InspectorPanel inspector = new InspectorPanel();
    
    // Свързване към Observer система
    c.simulation.snapshots().subscribe(map);
    c.simulation.snapshots().subscribe(timeline);
    c.simulation.snapshots().subscribe(inspector);
    
    // Медиатор за бутоните
    new MainWindowMediator(c.simulation, start, pause);
}
\end{lstlisting}

\section{Шаблони за проектиране (Design Patterns)}

\subsection{Създаващи шаблони (Creational Patterns)}

\subsubsection{Abstract Factory}

\textbf{Цел:} Предоставя интерфейс за създаване на семейства от свързани обекти.

\textbf{Имплементация:} Използва се за създаване на различни теми (Default, Sofia).

\begin{figure}[H]
\centering
\begin{tikzpicture}
\umlinterface[x=0, y=4]{CityPackFactory}{
    + palette() : Palette \\
    + cityName() : String \\
    + gridSize() : int
}

\umlclass[x=-3, y=0]{DefaultPackFactory}{}{
    + palette() : Palette \\
    + cityName() : String \\
    + gridSize() : int
}

\umlclass[x=3, y=0]{SofiaPackFactory}{}{
    + palette() : Palette \\
    + cityName() : String \\
    + gridSize() : int
}

\umlclass[x=6, y=4]{Palette}{
    + background : Color \\
    + grid : Color \\
    + vehicle : Color \\
    + accent : Color
}{}

\umlimpl{DefaultPackFactory}{CityPackFactory}
\umlimpl{SofiaPackFactory}{CityPackFactory}
\umlassoc{CityPackFactory}{Palette}
\end{tikzpicture}
\caption{UML диаграма на Abstract Factory pattern}
\end{figure}

\textbf{Предимства:}
\begin{itemize}
    \item Лесно добавяне на нови теми без промяна на съществуващия код
    \item Гарантира съвместимост между компонентите на дадена тема
    \item Поддържа плъгин архитектура чрез ServiceLoader
\end{itemize}

\subsubsection{Builder}

\textbf{Цел:} Позволява създаването на сложни обекти стъпка по стъпка.

\textbf{Имплементация:} \texttt{ScenarioBuilder} изгражда конфигурации на света.

\begin{figure}[H]
\centering
\begin{tikzpicture}
\umlclass[x=0, y=2]{ScenarioBuilder}{
    - world : World
}{
    + addVehicle(label, start, target) : ScenarioBuilder \\
    + build() : World
}

\umlclass[x=4, y=2]{World}{
    - vehicles : List<Vehicle> \\
    - clock : SimClock
}{
    + vehicles() : List<Vehicle> \\
    + clock() : SimClock
}

\umlclass[x=4, y=-1]{Vehicle}{
    - label : String \\
    - nodeIndex : int \\
    - state : VehicleState
}{
    + tick(clock, ct) : void
}

\umlassoc{ScenarioBuilder}{World}
\umlcompo{World}{Vehicle}
\end{tikzpicture}
\caption{UML диаграма на Builder pattern}
\end{figure}

\textbf{Предимства:}
\begin{itemize}
    \item Флуентен API за създаване на сценарии
    \item Възможност за различни конфигурации без множество конструктори
    \item Лесно тестване с различни настройки
\end{itemize}

\subsection{Структурни шаблони (Structural Patterns)}

\subsubsection{Facade}

\textbf{Цел:} Предоставя опростен интерфейс към сложна подсистема.

\textbf{Имплементация:} \texttt{SimulationFacade} скрива сложността на управлението на нишки.

\begin{figure}[H]
\centering
\begin{tikzpicture}
\umlclass[x=0, y=3]{SimulationFacade}{
    - scheduler : ScheduledExecutorService \\
    - pool : ExecutorService \\
    - engine : SimulationEngine \\
    - token : CancellationToken
}{
    + start() : void \\
    + pause() : void \\
    + snapshots() : Subject<WorldSnapshot>
}

\umlclass[x=6, y=4]{SimulationEngine}{
    - pool : ExecutorService \\
    - world : World
}{
    + step(ct) : WorldSnapshot
}

\umlclass[x=6, y=1]{Subject<WorldSnapshot>}{
    - observers : List<Observer>
}{
    + subscribe(o) : void \\
    + publish(value) : void
}

\umlassoc{SimulationFacade}{SimulationEngine}
\umlassoc{SimulationFacade}{Subject<WorldSnapshot>}
\end{tikzpicture}
\caption{UML диаграма на Facade pattern}
\end{figure}

\textbf{Предимства:}
\begin{itemize}
    \item Скрива сложността на паралелното изпълнение
    \item Предоставя прост интерфейс за клиентите
    \item Централизира управлението на ресурсите
\end{itemize}

\subsubsection{Decorator}

\textbf{Цел:} Добавя ново поведение към обекти динамично без промяна на тяхната структура.

\textbf{Имплементация:} \texttt{MeteredRouter} и \texttt{LoggingRouter} добавят функционалност към базовите маршрутизатори.

\begin{figure}[H]
\centering
\begin{tikzpicture}
\umlinterface[x=0, y=4]{Router}{
    + compute(req, ct) : Route
}

\umlclass[x=-4, y=1]{TimeOptimalRouter}{}{
    + compute(req, ct) : Route
}

\umlclass[x=0, y=1]{LoggingRouter}{
    - inner : Router \\
    - metrics : MetricsRegistry
}{
    + compute(req, ct) : Route
}

\umlclass[x=4, y=1]{MeteredRouter}{
    - inner : Router \\
    - metrics : MetricsRegistry
}{
    + compute(req, ct) : Route
}

\umlimpl{TimeOptimalRouter}{Router}
\umlimpl{LoggingRouter}{Router}
\umlimpl{MeteredRouter}{Router}
\umlassoc{LoggingRouter}{Router}
\umlassoc{MeteredRouter}{Router}
\end{tikzpicture}
\caption{UML диаграма на Decorator pattern}
\end{figure}

\textbf{Предимства:}
\begin{itemize}
    \item Гъвкаво композиране на функционалности
    \item Спазване на Open/Closed principle
    \item Възможност за многократно декориране
\end{itemize}

\subsubsection{Adapter}

\textbf{Цел:} Позволява на класове с несъвместими интерфейси да работят заедно.

\textbf{Имплементация:} \texttt{ThirdPartyRouterAdapter} адаптира външна библиотека за маршрутизиране.

\begin{figure}[H]
\centering
\begin{tikzpicture}
\umlinterface[x=0, y=3]{Router}{
    + compute(req, ct) : Route
}

\umlclass[x=3, y=0]{ThirdPartyPathfinder}{}{
    + findPath(src, dst) : List<Integer>
}

\umlclass[x=0, y=0]{ThirdPartyRouterAdapter}{
    - lib : ThirdPartyPathfinder
}{
    + compute(req, ct) : Route
}

\umlimpl{ThirdPartyRouterAdapter}{Router}
\umlassoc{ThirdPartyRouterAdapter}{ThirdPartyPathfinder}
\end{tikzpicture}
\caption{UML диаграма на Adapter pattern}
\end{figure}

\textbf{Предимства:}
\begin{itemize}
    \item Интеграция на външни библиотеки без промяна в тях
    \item Запазване на съществуващите интерфейси
    \item Изолиране на зависимостите към външни системи
\end{itemize}

\subsection{Поведенчески шаблони (Behavioral Patterns)}

\subsubsection{Observer}

\textbf{Цел:} Дефинира зависимост тип "един-към-много" между обекти.

\textbf{Имплементация:} \texttt{Subject} публикува snapshot-и към UI панелите.

\begin{figure}[H]
\centering
\begin{tikzpicture}
\umlclass[x=0, y=4]{Subject<T>}{
    - observers : List<Observer<T>>
}{
    + subscribe(o) : void \\
    + unsubscribe(o) : void \\
    + publish(value) : void
}

\umlinterface[x=5, y=4]{Observer<T>}{
    + onNext(value) : void
}

\umlclass[x=3, y=1]{MapPanel}{}{
    + onNext(snapshot) : void
}

\umlclass[x=5, y=1]{TimelinePanel}{}{
    + onNext(snapshot) : void
}

\umlclass[x=7, y=1]{InspectorPanel}{}{
    + onNext(snapshot) : void
}

\umlassoc{Subject<T>}{Observer<T>}
\umlimpl{MapPanel}{Observer<T>}
\umlimpl{TimelinePanel}{Observer<T>}
\umlimpl{InspectorPanel}{Observer<T>}
\end{tikzpicture}
\caption{UML диаграма на Observer pattern}
\end{figure}

\textbf{Предимства:}
\begin{itemize}
    \item Слабо свързване между издател и наблюдатели
    \item Динамично добавяне/премахване на наблюдатели
    \item Поддържа broadcast комуникация
\end{itemize}

\subsubsection{State}

\textbf{Цел:} Позволява на обект да променя поведението си при промяна на вътрешното състояние.

\textbf{Имплементация:} \texttt{Vehicle} преминава през различни състояния (Idle → EnRoute → Dwell).

\begin{figure}[H]
\centering
\begin{tikzpicture}
\umlinterface[x=4, y=4]{VehicleState}{
    + onTick(ctx, clock, ct) : void \\
    + name() : String
}

\umlclass[x=0, y=2]{Vehicle}{
    - state : VehicleState \\
    - nodeIndex : int \\
    - targetIndex : int \\
    - label : String
}{
    + tick(clock, ct) : void \\
    + setState(s) : void \\
    + stateName() : String
}

\umlclass[x=2, y=0]{Idle}{}{
    + onTick(ctx, clock, ct) : void \\
    + name() : String
}

\umlclass[x=4, y=0]{EnRoute}{}{
    + onTick(ctx, clock, ct) : void \\
    + name() : String
}

\umlclass[x=6, y=0]{Dwell}{
    - remaining : int
}{
    + onTick(ctx, clock, ct) : void \\
    + name() : String
}

\umlassoc{Vehicle}{VehicleState}
\umlimpl{Idle}{VehicleState}
\umlimpl{EnRoute}{VehicleState}
\umlimpl{Dwell}{VehicleState}
\end{tikzpicture}
\caption{UML диаграма на State pattern}
\end{figure}

\textbf{Логика на преходите:}
\begin{itemize}
    \item \textbf{Idle → EnRoute:} При започване на симулацията
    \item \textbf{EnRoute → Dwell:} При достигане на дестинация
    \item \textbf{Dwell → EnRoute:} След изтичане на времето за престой
\end{itemize}

\subsubsection{Strategy}

\textbf{Цел:} Дефинира семейство от алгоритми и ги прави взаимозаменяеми.

\textbf{Имплементация:} \texttt{Router} интерфейсът позволява различни алгоритми за маршрутизиране.

\begin{figure}[H]
\centering
\begin{tikzpicture}
\umlinterface[x=0, y=3]{Router}{
    + compute(req, ct) : Route
}

\umlclass[x=-3, y=0]{TimeOptimalRouter}{}{
    + compute(req, ct) : Route
}

\umlclass[x=3, y=0]{ThirdPartyRouterAdapter}{
    - lib : ThirdPartyPathfinder
}{
    + compute(req, ct) : Route
}

\umlclass[x=0, y=0]{CustomRouter}{}{
    + compute(req, ct) : Route
}

\umlimpl{TimeOptimalRouter}{Router}
\umlimpl{ThirdPartyRouterAdapter}{Router}
\umlimpl{CustomRouter}{Router}
\end{tikzpicture}
\caption{UML диаграма на Strategy pattern}
\end{figure}

\textbf{Стратегии:}
\begin{itemize}
    \item \textbf{TimeOptimalRouter:} Най-кратко време за пътуване
    \item \textbf{ThirdPartyRouterAdapter:} Използва външна библиотека
    \item \textbf{CustomRouter:} Потребителски алгоритми (възможност за разширение)
\end{itemize}

\subsubsection{Command + Memento}

\textbf{Цел:} Капсулира заявка като обект и поддържа undo/redo операции.

\textbf{Имплементация:} \texttt{CommandStack} осигурява scaffold за бъдещи editing функции.

\begin{figure}[H]
\centering
\begin{tikzpicture}
\umlinterface[x=0, y=3]{Command}{
    + execute() : void \\
    + undo() : void
}

\umlclass[x=4, y=3]{CommandStack}{
    - undo : Deque<Command> \\
    - redo : Deque<Command>
}{
    + apply(c) : void \\
    + undo() : void \\
    + redo() : void
}

\umlclass[x=0, y=0]{ConcreteCommand}{
    - receiver : Object \\
    - state : Memento
}{
    + execute() : void \\
    + undo() : void
}

\umlassoc{CommandStack}{Command}
\umlimpl{ConcreteCommand}{Command}
\end{tikzpicture}
\caption{UML диаграма на Command + Memento pattern}
\end{figure}

\textbf{Възможности за разширение:}
\begin{itemize}
    \item Команди за добавяне/премахване на превозни средства
    \item Команди за промяна на маршрути
    \item Команди за конфигуриране на симулацията
\end{itemize}

\subsubsection{Mediator}

\textbf{Цел:} Определя как набор от обекти взаимодействат помежду си.

\textbf{Имплементация:} \texttt{MainWindowMediator} координира взаимодействието между бутоните и симулацията.

\begin{figure}[H]
\centering
\begin{tikzpicture}
\umlclass[x=0, y=2]{MainWindowMediator}{
    - sim : SimulationFacade \\
    - startBtn : JButton \\
    - pauseBtn : JButton
}{
    + wire() : void
}

\umlclass[x=-3, y=-1]{JButton}{}{
    + addActionListener(l) : void
}

\umlclass[x=3, y=-1]{SimulationFacade}{}{
    + start() : void \\
    + pause() : void
}

\umlassoc{MainWindowMediator}{JButton}
\umlassoc{MainWindowMediator}{SimulationFacade}
\end{tikzpicture}
\caption{UML диаграма на Mediator pattern}
\end{figure}

\textbf{Координирани действия:}
\begin{itemize}
    \item Start бутонът стартира симулацията
    \item Pause бутонът паузира симулацията
    \item Обработка на грешки чрез диалогови прозорци
\end{itemize}

\section{Конкурентност и паралелизъм}

\subsection{Модел на нишки}

MetroSim използва софистициран модел на паралелност за да осигури:
\begin{itemize}
    \item \textbf{Отзивчив потребителски интерфейс}
    \item \textbf{Ефективна обработка на превозните средства}
    \item \textbf{Безопасна синхронизация}
\end{itemize}

\subsection{Архитектура на нишките}

\begin{figure}[H]
\centering
\begin{tikzpicture}
% EDT Thread
\draw[fill=lightblue, rounded corners] (0,6) rectangle (3,8);
\node[align=center] at (1.5,7) {\textbf{EDT Thread} \\ UI Updates \\ User Interaction};

% Scheduler Thread
\draw[fill=lightgreen, rounded corners] (4,6) rectangle (7,8);
\node[align=center] at (5.5,7) {\textbf{Scheduler} \\ Periodic Ticks \\ 50ms intervals};

% Worker Pool
\draw[fill=lightyellow, rounded corners] (8,6) rectangle (11,8);
\node[align=center] at (9.5,7) {\textbf{Worker Pool} \\ Vehicle Updates \\ Parallel Execution};

% Data Flow
\draw[->, thick] (5.5,6) -- (5.5,4.5) -- (1.5,4.5) -- (1.5,6);
\node[above] at (3.5,4.5) {WorldSnapshot};

\draw[->, thick] (5.5,6) -- (5.5,3.5) -- (9.5,3.5) -- (9.5,6);
\node[below] at (7.5,3.5) {Parallel Tasks};

% Synchronization
\draw[fill=lightcoral, rounded corners] (4,1) rectangle (7,3);
\node[align=center] at (5.5,2) {\textbf{CancellationToken} \\ Thread-safe \\ Coordination};

\draw[dashed] (5.5,3) -- (5.5,6);
\draw[dashed] (5.5,3) -- (9.5,6);
\end{tikzpicture}
\caption{Архитектура на нишките в MetroSim}
\end{figure}

\subsection{Синхронизационни механизми}

\subsubsection{CancellationToken}

Осигурява кооперативно прекъсване на операции:

\begin{lstlisting}[caption=CancellationToken имплементация]
public final class CancellationToken {
    private final AtomicBoolean cancelled = new AtomicBoolean(false);
    
    public boolean isCancelled() {
        return cancelled.get();
    }
    
    public void cancel() {
        cancelled.set(true);
    }
}
\end{lstlisting}

\subsubsection{Неизменяеми обекти (Immutable Objects)}

Всички данни, които се предават между нишки, са неизменяеми:

\begin{itemize}
    \item \texttt{WorldSnapshot} – състояние на света
    \item \texttt{VehicleSnapshot} – състояние на превозно средство
    \item \texttt{Route} – маршрут между две точки
\end{itemize}

\subsection{Предотвратяване на състезателни условия}

\begin{enumerate}
    \item \textbf{Copy-on-read:} Снимките се правят преди публикуване
    \item \textbf{Атомични операции:} Използват се за ъпдейт на състоянието
    \item \textbf{Synchronized блокове:} За критичните секции в SimulationFacade
    \item \textbf{Volatile променливи:} За видимост между нишки
\end{enumerate}

\section{Метрики и наблюдение}

\subsection{MetricsRegistry система}

Системата за метрики позволява проследяване на производителността:

\begin{lstlisting}[caption=MetricsRegistry интерфейс]
public interface MetricsRegistry {
    void counter(String name, long delta);
    void recordTime(String name, long nanos);
    long getCounter(String name);
    long getTimeSum(String name);
}
\end{lstlisting}

\subsection{Типове метрики}

\begin{table}[H]
\centering
\begin{tabular}{|p{4cm}|p{9cm}|}
\hline
\textbf{Тип метрика} & \textbf{Описание} \\
\hline
routing.count & Общ брой извършени маршрутизации \\
\hline
routing.total.ns & Общо време за маршрутизация в наносекунди \\
\hline
routing.logged & Брой логнати маршрутизации \\
\hline
simulation.ticks & Общ брой симулационни tick-ове \\
\hline
\end{tabular}
\caption{Поддържани метрики в системата}
\end{table}

\subsection{Декоратори за метрики}

Метриките се събират чрез Decorator pattern:

\begin{lstlisting}[caption=MeteredRouter декоратор]
public Route compute(RouteRequest req, CancellationToken ct) {
    long t0 = System.nanoTime();
    try {
        return inner.compute(req, ct);
    } finally {
        metrics.recordTime("routing.total.ns", System.nanoTime() - t0);
        metrics.counter("routing.count", 1);
    }
}
\end{lstlisting}

\section{Система за плъгини и разширяемост}

\subsection{ServiceLoader архитектура}

MetroSim използва Java ServiceLoader за откриване на плъгини по време на изпълнение:

\begin{lstlisting}[caption=Плъгин discovery в Bootstrap]
private static CityPackFactory resolveCityPack() {
    String want = System.getProperty("metrosim.cityPack", "");
    CityPackFactory fallback = new DefaultPackFactory();
    var loader = ServiceLoader.load(CityPackFactory.class);
    
    for (CityPackFactory f : loader) {
        String name = f.cityName().toLowerCase();
        if (!want.isEmpty() && name.contains(want)) {
            return f;
        }
    }
    return fallback;
}
\end{lstlisting}

\subsection{Конфигурационен файл}

Плъгините се регистрират в \texttt{META-INF/services/metrosim.themes.CityPackFactory}:

\begin{lstlisting}[caption=Регистрация на плъгини]
metrosim.themes.DefaultPackFactory
metrosim.themes.SofiaPackFactory
\end{lstlisting}

\subsection{Примерни плъгини}

\subsubsection{Sofia Pack Factory}

\begin{lstlisting}[caption=SofiaPackFactory имплементация]
public final class SofiaPackFactory implements CityPackFactory {
    public Palette palette() {
        return new Palette(
            new Color(0x0f1220), // background
            new Color(0x2d355c), // grid
            new Color(0x58d68d), // vehicle
            new Color(0xf7dc6f)  // accent
        );
    }
    
    public String cityName() { return "Sofia"; }
    public int gridSize() { return 22; }
}
\end{lstlisting}

\subsection{Точки за разширение}

Системата предоставя следните точки за разширение:

\begin{itemize}
    \item \textbf{CityPackFactory} – нови визуalni теми
    \item \textbf{Router} – нови алгоритми за маршрутизиране
    \item \textbf{VehicleState} – нови поведения на превозните средства
    \item \textbf{Command} – нови команди за редактиране
    \item \textbf{Observer} – нови типове панели за визуализация
\end{itemize}

\section{Тестване и качество на кода}

\subsection{Архитектурни принципи}

Проектът следва основни принципи за качествен софтуер:

\begin{itemize}
    \item \textbf{SOLID принципи} – всеки клас има ясна отговорност
    \item \textbf{DRY (Don't Repeat Yourself)} – избягване на дублиране на код
    \item \textbf{KISS (Keep It Simple, Stupid)} – простота в дизайна
    \item \textbf{Separation of Concerns} – ясно разделение на отговорностите
\end{itemize}

\subsection{Кодови конвенции}

Проектът използва:
\begin{itemize}
    \item \textbf{Google Java Format} за форматиране на кода
    \item \textbf{Spotless} за автоматична корекция на стила
    \item \textbf{Gradle} за управление на build процеса
\end{itemize}

\subsection{Паралелно тестване}

Архитектурата позволява лесно unit тестване:

\begin{itemize}
    \item \textbf{Injection на зависимости} чрез конструктори
    \item \textbf{Интерфейси} за всички основни компоненти
    \item \textbf{Неизменяеми обекти} за предвидимо поведение
    \item \textbf{CancellationToken} за контролирано прекъсване
\end{itemize}

\section{Инструкции за стартиране и конфигуриране}

\subsection{Системни изисквания}

\begin{itemize}
    \item Java Development Kit (JDK) 17 или по-нова версия
    \item Операционна система с поддръжка на Java Swing
    \item Минимум 512MB RAM за изпълнение
\end{itemize}

\subsection{Компилиране и стартиране}

\subsubsection{Основни команди}

\begin{lstlisting}[language=bash, caption=Основни Gradle команди]
# Компилиране на проекта
./gradlew build

# Стартиране на приложението
./gradlew run

# Стартиране с конкретна тема
./gradlew run -Dmetrosim.cityPack=sofia

# Почистване на build артефакти
./gradlew clean
\end{lstlisting}

\subsubsection{Конфигурационни опции}

\begin{table}[H]
\centering
\begin{tabular}{|p{4cm}|p{9cm}|}
\hline
\textbf{Системно свойство} & \textbf{Описание} \\
\hline
metrosim.cityPack & Избор на градска тема (default, sofia) \\
\hline
java.awt.headless & Задава се на false за GUI режим \\
\hline
\end{tabular}
\caption{Конфигурационни параметри}
\end{table}

\subsection{Структура на проекта}

\begin{lstlisting}[caption=Структура на директориите]
metrosim/
├── src/main/java/metrosim/     # Основен код
│   ├── app/                    # Стартиране и bootstrap
│   ├── domain/                 # Бизнес логика
│   ├── infrastructure/         # Инфраструктура
│   ├── services/              # Услуги
│   ├── themes/                # Теми
│   ├── ui/                    # Потребителски интерфейс
│   └── events/                # Събитийна система
├── src/main/resources/         # Ресурси
│   ├── META-INF/services/     # ServiceLoader конфигурация
│   └── scenarios/             # Сценарии за симулация
├── build.gradle               # Build конфигурация
└── documentation/             # Документация
\end{lstlisting}

\section{Бъдещо развитие и разширения}

\subsection{Планирани функции}

\begin{itemize}
    \item \textbf{Редактор на сценарии} – графичен интерфейс за създаване на сценарии
    \item \textbf{Записване и възпроизвеждане} – запазване на симулации
    \item \textbf{Статистики в реално време} – графично представяне на метриките
    \item \textbf{Мрежови режим} – симулация на множество клиенти
    \item \textbf{3D визуализация} – по-реалистично представяне
\end{itemize}

\subsection{Възможни разширения}

\subsubsection{Нови типове превозни средства}
\begin{itemize}
    \item Автобуси с различни капацитети
    \item Трамваи с фиксирани релси
    \item Товарни превозни средства
\end{itemize}

\subsubsection{Сложни маршрути}
\begin{itemize}
    \item Условно маршрутизиране
    \item Динамично пренасочване
    \item Оптимизация по няколко критерия
\end{itemize}

\subsubsection{Интерактивност}
\begin{itemize}
    \item Добавяне на превозни средства по време на симулация
    \item Промяна на маршрути в реално време
    \item Симулиране на инциденти и закъснения
\end{itemize}

\section{Заключение}

\subsection{Постигнати цели}

MetroSim успешно демонстрира:

\begin{itemize}
    \item \textbf{Практическо приложение} на девет различни Design Patterns
    \item \textbf{Безопасна паралелност} в Java Swing приложения
    \item \textbf{Модулна архитектура} с ясно разделение на отговорностите
    \item \textbf{Разширяема система} за плъгини и теми
    \item \textbf{Реално време симулация} с отзивчив потребителски интерфейс
\end{itemize}

\subsection{Образователна стойност}

Проектът предоставя отлична основа за обучение по:

\begin{itemize}
    \item \textbf{Object-Oriented Design} – правилно използване на наследяване и полиморфизъм
    \item \textbf{Concurrent Programming} – безопасна работа с множество нишки
    \item \textbf{Software Architecture} – изграждане на сложни, поддържаеми системи
    \item \textbf{Design Patterns} – практическо приложение в реален проект
\end{itemize}

\subsection{Техническо качество}

Кодовата база поддържа високи стандарти за качество:

\begin{itemize}
    \item \textbf{Четим и добре документиран код}
    \item \textbf{Consistent code style} чрез автоматично форматиране
    \item \textbf{Модулна структура} за лесно разширение
    \item \textbf{Ефективно използване на паметта} чрез immutable обекти
    \item \textbf{Thread-safe архитектура} без състезателни условия
\end{itemize}

\subsection{Практическа приложимост}

Принципите и шаблоните, демонстрирани в MetroSim, могат да се приложат в:

\begin{itemize}
    \item \textbf{Реални симулационни системи}
    \item \textbf{Игрови двигатели}
    \item \textbf{Финансови системи за трейдинг}
    \item \textbf{IoT и embedded системи}
    \item \textbf{Микросервизна архитектура}
\end{itemize}

MetroSim представлява солиден пример за това как правилното прилагане на Design Patterns и принципите за добра софтуерна архитектура могат да доведат до създаването на мащабируеми, поддържаеми и разширяеми системи.

\newpage

\section*{Приложения}

\subsection*{A. Пълен UML клас диаграм на системата}

\begin{figure}[H]
\centering
\begin{tikzpicture}[scale=0.8]
% Domain Layer
\umlclass[x=0, y=8]{World}{
    - vehicles : List<Vehicle> \\
    - clock : SimClock
}{
    + vehicles() : List<Vehicle> \\
    + clock() : SimClock \\
    + advanceClock() : void \\
    + snapshot(tickId) : WorldSnapshot
}

\umlclass[x=4, y=8]{Vehicle}{
    - state : VehicleState \\
    - nodeIndex : int \\
    - targetIndex : int \\
    - label : String
}{
    + tick(clock, ct) : void \\
    + setState(s) : void \\
    + stateName() : String
}

% Infrastructure Layer
\umlclass[x=0, y=4]{SimulationEngine}{
    - pool : ExecutorService \\
    - world : World \\
    - tick : AtomicLong
}{
    + step(ct) : WorldSnapshot
}

\umlclass[x=4, y=4]{SimulationFacade}{
    - scheduler : ScheduledExecutorService \\
    - engine : SimulationEngine \\
    - snapshots : Subject<WorldSnapshot>
}{
    + start() : void \\
    + pause() : void \\
    + snapshots() : Subject<WorldSnapshot>
}

% UI Layer
\umlclass[x=8, y=8]{MainWindow}{
    - components : Bootstrap.Components
}{
    + MainWindow(c) : void
}

\umlclass[x=8, y=4]{MapPanel}{
    - snapshot : WorldSnapshot \\
    - packs : CityPackFactory
}{
    + onNext(s) : void \\
    + paintComponent(g) : void
}

% Services Layer
\umlinterface[x=0, y=0]{Router}{
    + compute(req, ct) : Route
}

\umlclass[x=4, y=0]{MeteredRouter}{
    - inner : Router \\
    - metrics : MetricsRegistry
}{
    + compute(req, ct) : Route
}

% Associations
\umlcompo{World}{Vehicle}
\umlassoc{SimulationEngine}{World}
\umlassoc{SimulationFacade}{SimulationEngine}
\umlassoc{MainWindow}{MapPanel}
\umlimpl{MeteredRouter}{Router}
\end{tikzpicture}
\caption{Опростен UML клас диаграм на цялата система}
\end{figure}

\subsection*{B. Последователност диаграм на симулационния цикъл}

\begin{figure}[H]
\centering
\begin{tikzpicture}
\umlactor{User}
\umlentity[x=3]{MainWindow}
\umlentity[x=6]{SimulationFacade}
\umlentity[x=9]{SimulationEngine}
\umlentity[x=12]{Vehicle}

\umlcall[dt=5]{User}{MainWindow}{start button click}
\umlcall[dt=5]{MainWindow}{SimulationFacade}{start()}
\umlcall[dt=5]{SimulationFacade}{SimulationEngine}{step(token)}
\umlcall[dt=5]{SimulationEngine}{Vehicle}{tick(clock, token)}
\umlreturn[dt=5]{Vehicle}{SimulationEngine}{}
\umlreturn[dt=5]{SimulationEngine}{SimulationFacade}{WorldSnapshot}
\umlcall[dt=5]{SimulationFacade}{MainWindow}{publish(snapshot)}
\end{tikzpicture}
\caption{Последователност диаграм на симулационния цикъл}
\end{figure}

\subsection*{C. Справочник на ключовите класове}

\begin{table}[H]
\centering
\small
\begin{tabular}{|p{3cm}|p{4cm}|p{7cm}|}
\hline
\textbf{Клас} & \textbf{Пакет} & \textbf{Отговорност} \\
\hline
Bootstrap & app & Инициализация и конфигуриране на компонентите \\
\hline
MetroSimApp & app & Главна точка на стартиране \\
\hline
World & domain & Централно състояние на симулацията \\
\hline
Vehicle & domain & Превозно средство с вътрешни състояния \\
\hline
SimulationEngine & infrastructure & Ядро на симулацията с паралелна обработка \\
\hline
SimulationFacade & infrastructure & Опростен интерфейс за управление на симулацията \\
\hline
Router & services.routing & Стратегия за изчисляване на маршрути \\
\hline
CityPackFactory & themes & Фабрика за създаване на визуални теми \\
\hline
Subject & events & Събитийна система за Observer pattern \\
\hline
MainWindow & ui & Главен прозорец на приложението \\
\hline
CommandStack & ui.commands & Система за undo/redo операции \\
\hline
\end{tabular}
\caption{Справочник на ключовите класове}
\end{table}

\end{document}